% Part: model-theory
% Chapter: models-of-arithmetic
% Section: standard-models

\documentclass[../../../include/open-logic-section]{subfiles}

\begin{document}

\olfileid{mod}{mar}{stm}
\olsection{نماذج الحساب القياسية}

إنما وُضعت لغة الحساب~$\Lang{L_A}$ للكلام في الأعداد الطبيعية.
ولهذا اختص «النموذج» القياسي $\Struct{N}$ بمنزلة، إذ هو المقصود
أولًا بالحديث. غير أن نظرنا في المنطق كثيرًا ما يتعلق بالخواص
البنيوية وحدها؛ فإذا تشاكلت اثنتان من ال!!{structure}s
اتفقتا في تلك الخواص. فلنا إذن أن نوسّع الاصطلاح، ونسمي كل !!{structure} تتشاكل
مع~$\Struct{N}$ «قياسية».

\begin{defn}
تكون !!{structure} للغة $\Lang{L_A}$ \emph{قياسية} إذا كانت متشاكلة
مع~$\Struct{N}$.
\end{defn}

\begin{prop}
\ollabel{prop:standard-domain} إذا كانت !!a{structure}~$\Struct{M}$
قياسية، فإن مجالها مجموعة قيم الأعداد القياسية، أي:
\[
\Domain{M} = \Setabs{\Value{\num{n}}{M}}{n \in \Nat}
\]
\end{prop}

\begin{proof}
لدينا أولًا $\Value{\num{n}}{M} \in \Domain{M}$ لكل عدد قياسي.
ويبقى إثبات أن كل $x \in \Domain{M}$ يساوي $\Value{\num{n}}{M}$
لبعض~$n$. وبما أن $\Struct{M}$ قياسية، فهي متشاكلة مع~$\Struct{N}$؛
فلنأخذ تشاكلًا $g\colon \Nat \to \Domain{M}$. عندئذ
$g(n) = g(\Value{\num{n}}{N}) = \Value{\num{n}}{M}$.
ولكل $x \in \Domain{M}$ يوجد $n \in \Nat$ يحقق $g(n)=x$،
لأن $g$~!!{surjective}. فتم الاحتواء الآخر.
\end{proof}

\begin{explain}
فالبنية~$\Struct{M}$ للغة $\Lang{L_A}$، متى كانت قياسية،
أمكن تسمية كل عنصر في !!{domain}ها بالحدود العددية القياسية
$\num{0}$، $\num{1}$، $\num{2}$، \dots؛ أي بالحدود
$\Obj{0}$، $\Obj{0}'$، $\Obj{0}''$، وهكذا. وليس معناه أن
!!{element}s~$\Domain{M}$ \emph{هي} الأعداد أنفسها؛ وإنما نعيّن
تلك العناصر بالطريقة التي نعيّن بها الأعداد في~$\Domain{N}$.
\end{explain}

\begin{prob}
بيّن أن عكس \olref[mod][mar][stm]{prop:standard-domain} كاذب؛ أي أعط
مثالًا لـ!!a{structure}~$\Struct{M}$ تحقق
$\Domain{M} = \Setabs{\Value{\num{n}}{M}}{n \in \Nat}$ ولا تكون
متشاكلة مع~$\Struct{N}$.
\end{prob}

\begin{prop}
\ollabel{prop:thq-standard}
إذا كان $\Sat{M}{\Th{Q}}$ و
$\Domain{M} = \Setabs{\Value{\num{n}}{M}}{n \in \Nat}$، فإن
$\Struct{M}$ قياسية.
\end{prop}

\begin{proof}
المطلوب تشاكل $\Struct{M}$ مع~$\Struct{N}$. فلنعرّف
$g\colon \Nat \to \Domain{M}$ بوضع
$g(n)=\Value{\num{n}}{M}$. الفرض يقضي بأن $g$~!!{surjective}.
وهي كذلك !!{injective}: فإن $n \neq m$ يستلزم
$\Th{Q} \Proves \eq/[\num{n}][\num{m}]$. ولما كان
$\Sat{M}{\Th{Q}}$، لزم $\Sat{M}{\eq/[\num{n}][\num{m}]}$
كلما كان $n \neq m$. فإذا كان $n \neq m$، كان
$\Value{\num{n}}{M} \neq \Value{\num{m}}{M}$، أي $g(n) \neq g(m)$.

وبقي التحقق من حفظ $g$ لتأويلات الرموز، فنفصلها واحدًا واحدًا:
\begin{enumerate}
\item لدينا $g(\Assign{\Obj{0}}{N})=g(0)$، لأن
  $\Assign{\Obj{0}}{N}=0$. وبتعريف~$g$، يكون
  $g(0)=\Value{\num{0}}{M}$. لكن $\num{0}$ هو $\Obj{0}$ نفسه، وقيمة
  الحد الذي يكون !!a{constant} هي ما تسنده الـ!!{structure} إلى ذلك
  الـ!!{constant}، أي
  $\Value{\Obj{0}}{M}=\Assign{\Obj{0}}{M}$. وهكذا نحصل، كما يلزم، على
  $g(\Assign{\Obj{0}}{N})=\Assign{\Obj{0}}{M}$.
\item $g(\Assign{\prime}{N}(n))=g(n+1)$، لأن $\prime$ في
  $\Struct{N}$ دالة الخلف على~$\Nat$. ثم
  $g(n+1)=\Value{\num{n+1}}{M}$ بتعريف~$g$. لكن $\num{n+1}$ هو الحد
  نفسه $\num{n}'$، ومن ثم
  $\Value{\num{n+1}}{M}=\Value{\num{n}'}{M}$. وبتعريف دالة القيمة،
  يساوي هذا $\Assign{\prime}{M}(\Value{\num{n}}{M})$. وبما أن
  $\Value{\num{n}}{M}=g(n)$، نحصل على
  $g(\Assign{\prime}{N}(n))=\Assign{\prime}{M}(g(n))$.
\item $g(\Assign{+}{N}(n,m))=g(n+m)$، لأن $+$ في $\Struct{N}$ دالة
  الجمع على~$\Nat$. ثم $g(n+m)=\Value{\num{n+m}}{M}$ بتعريف~$g$.
  لكن $\Th{Q} \Proves \num{n+m} = (\num{n} + \num{m})$، ومن ثم
  $\Value{\num{n+m}}{M}=\Value{\num{n}+\num{m}}{M}$. وبتعريف دالة
  القيمة يساوي هذا
  $\Assign{+}{M}(\Value{\num{n}}{M},\Value{\num{m}}{M})$. وبما أن
  $\Value{\num{n}}{M}=g(n)$ و$\Value{\num{m}}{M}=g(m)$، نحصل على
  $g(\Assign{+}{N}(n,m))=\Assign{+}{M}(g(n),g(m))$.
\item $g(\Assign{\times}{N}(n, m)) = \Assign{\times}{M}(g(n), g(m))$:
  تمرين.
\item $\tuple{n,m} \in \Assign{<}{N}$ إذا وفقط إذا كان $n<m$. فإذا
  كان $n<m$، كان $\Th{Q} \Proves \num{n} < \num{m}$، وكذلك
  $\Sat{M}{\num{n} < \num{m}}$. ومن ثم
  $\tuple{\Value{\num{n}}{M},\Value{\num{m}}{M}} \in \Assign{<}{M}$،
  أي $\tuple{g(n),g(m)} \in \Assign{<}{M}$. وإذا كان $n \not< m$، كان
  $\Th{Q} \Proves \lnot \num{n} < \num{m}$، وبالتالي
  $\Sat/{M}{\num{n} < \num{m}}$. وهكذا، كما سبق،
  $\tuple{g(n),g(m)} \notin \Assign{<}{M}$. وبضم الحالتين نحصل على:
  $\tuple{n,m} \in \Assign{<}{N}$ إذا وفقط إذا كان
  $\tuple{g(n),g(m)} \in \Assign{<}{M}$.
\end{enumerate}
\end{proof}

\begin{explain}
الدالة~$g$ أقرب ما نعرّف به تطبيقًا من $\Nat$ إلى مجال
!!{structure}~$\Struct{M}$ للغة $\Lang{L_A}$. فإن كل
$\Struct{M}$ تشتمل على !!{element}s تسميها الحدود
$\num{0}$ و$\num{1}$ و$\num{2}$ وهكذا. فإذا وافقت $\Struct{M}$
ما يصدق عن الـ$\num{n}$ في~$\Struct{N}$ من العبارات الأساسية
المذكورة، وكانت $g$ تقابلًا، لم يبعد أن تكون $g$ تشاكلًا.
بل هي التشاكل الوحيد إذا لم يكن في $\Domain{M}$ من
!!{element}s إلا ما تسميه الـ$\num{n}$.
\end{explain}

\begin{prop}
\ollabel{prop:thq-unique-iso} إذا كانت $\Struct{M}$ قياسية، كانت $g$ من
برهان \olref{prop:thq-standard} التشاكل الوحيد من $\Struct{N}$ إلى~$\Struct{M}$.
\end{prop}

\begin{proof}
ليكن $h\colon \Nat \to \Domain{M}$ تشاكلًا بين $\Struct{N}$
و$\Struct{M}$. نستقرئ على~$n$ لإثبات المساواة $g=h$.
عند $n=0$ لدينا $g(0)=\Assign{\Obj{0}}{M}$ بتعريف~$g$.
ولأن $h$ تشاكل، فـ
$h(0)=h(\Assign{\Obj{0}}{N})=\Assign{\Obj{0}}{M}$؛
ومن ثم $g(0)=h(0)$.

وأما في خطوة $n+1$، فالحساب الآتي يتم الاستقراء:
\begin{align*}
  g(n+1) & = \Value{\num{n+1}}{M} \text{ بتعريف~$g$}\\
  & = \Value{\num{n}'}{M} \text{ لأن $\num{n+1}\ident \num{n}'$}\\
  & = \Assign{\prime}{M}(\Value{\num{n}}{M})
    \text{ بتعريف $\Value{t'}{M}$}\\
  & = \Assign{\prime}{M}(g(n)) \text{ بتعريف~$g$}\\
  & = \Assign{\prime}{M}(h(n)) \text{ بفرض الاستقراء}\\
  & = h(\Assign{\prime}{N}(n)) \text{ لأن $h$ تشاكل}\\
  & = h(n+1)
\end{align*}
\end{proof}

\begin{explain}
كل مجموعة~$M$~!!{denumerable} تقبل !!a{bijection} بين
$\Nat$ و$M$؛ فكل مجموعة~$M$ كهذه تصلح !!{domain} لنموذج
قياسي $\Struct{M}$. فإذا اخترنا العنصر $z \in M$ ودالة
مناسبة~$s$ لتأويل $\Assign{\Obj{0}}{M}$ و$\Assign{\prime}{M}$،
تعيّن تأويل $+$ و$\times$ و$<$ بذلك الاختيار.
ومن الشروط اللازمة لتأويل $\Assign{\prime}{M}$ أن تكون الدالة
$s\colon M \to M \setminus \{z\}$~!!{injective} و!!{surjective}.
أما مدى~$s$ فلا يدخل فيه~$z$، وإلا كذبت
$\lforall[x][\eq/[\Obj 0][x']]$. وهذه !!{sentence} تصدق
في~$\Struct{N}$، فلا بد أن تصدق في $\Struct{M}$.
وأما $s$ فلا بد أن تكون !!{injective}، لأن دالة الخلف
$\Assign{\prime}{N}$ في $\Struct{N}$ كذلك؛ وكون
$\Assign{\prime}{N}$~!!{injective} تعبّر عنه !!a{sentence}
تصدق في $\Struct{N}$. ولا بد أيضًا أن تكون !!{surjective}؛
وإلا وجد $x \in M \setminus \{z\}$ خارج مدى~$s$، فتكذب
الـ!!{sentence}
$\lforall[x][(\eq[x][\Obj 0] \lor \lexists[y][\eq[y'][x]])]$
في~$\Struct{M}$، مع صدقها في~$\Struct{N}$.
\end{explain}


\end{document}
